% TOP_PROBLEM: {"id": 30006786, "title": "Parity conjecture for elliptic curves", "category_id": 1, "category": {"id": 1, "name": "number_theory", "display_name": "Number Theory", "description": "Properties of integers, prime numbers, Diophantine equations.", "slug": "number-theory", "order_index": 1, "created_at": "2026-07-31T15:26:25.670Z"}, "status": "open"}
% FIRST_POSTED: 2026-09-25
% ENTRY_KIND: statement_only
% !TeX program = lualatex
% Definition and sources only; this file is not an LLM solution attempt.
\documentclass[11pt]{article}
\usepackage[margin=25mm]{geometry}
\usepackage{fontspec}
\setmainfont{Latin Modern Roman}
\usepackage{amsmath,amssymb,mathtools}
\usepackage{unicode-math}
\setmathfont{Latin Modern Math}
\usepackage{xurl,hyperref}
\hypersetup{colorlinks=true,urlcolor=blue}
\setcounter{secnumdepth}{0}
\setlength{\parindent}{0pt}
\setlength{\parskip}{5pt}
\setlength{\emergencystretch}{3em}
\begin{document}
\section*{\texorpdfstring{Parity conjecture for elliptic curves}{Parity conjecture for elliptic curves}}\label{30006786}
\subsection{Definitions and mathematical statement}
By the Mordell--Weil rank of $E/K$ mean the integer $r$ in
$E(K)\cong E(K)_{\mathrm{tors}}\oplus{\mathbb{Z}}^r$. The global root number is
$w(E/K)=\prod_v w(E/K_v)\in\{\pm1\}$, the product of normalized local epsilon-factor signs; almost every finite-place factor is one and each archimedean factor is $-1$. The target is
\[
 w(E/K)=(-1)^r
\]
for every elliptic curve over every number field. Root numbers are defined locally without presuming a global analytic continuation theorem. The claim concerns the algebraic rank, not merely parity of a conjectural analytic order of vanishing or a Selmer-group rank.
\par\smallskip\textit{Formulation and concept references:} \href{https://arxiv.org/pdf/1207.0431}{[S1]}.

\subsection{Short English statement}
Does the product of an elliptic curve's local root-number signs always determine whether its rational-point rank over the number field is even or odd?
\subsection{Sources}
\begin{itemize}
\item[S1]Kęstutis Česnavičius, The p-parity conjecture for elliptic curves with a p-isogeny; arXiv1207.0431v3 April8,2014; PDF internal date May22,2018.
\url{https://arxiv.org/pdf/1207.0431}.
\textit{Formulation source.}
\end{itemize}
\end{document}
